\newpage
\section*{第26章：有界量化}
\addcontentsline{toc}{section}{第26章：有界量化}

\begin{taplsolutionentrybox}
\phantomsection\label{sol:translator:26.2.1}
\textbf{26.2.1 解答}

S-Arrow 要求参数逆变、结果协变。目标
\(B\to Y<:X\to B\) 因而分解为
\[
\Gamma\vdash X<:B
\qquad\text{和}\qquad
\Gamma\vdash Y<:B.
\]
第一项由上下文 \(X<:B\) 和 S-TVar 得到；第二项先由 \(Y<:X\)、
\(X<:B\) 各用一次 S-TVar，再用 S-Trans 得到。最后应用 S-Arrow。
\taplsolutionbacklink{ex:26.2.1}{返回练习26.2.1}
\end{taplsolutionentrybox}

\begin{taplsolutionentrybox}
\phantomsection\label{sol:translator:26.2.2}
\textbf{26.2.2 解答}

取 \(A<:B\)。完全规则可推出
\[
\forall X<:B.X
<:
\forall X<:A.B,
\]
因为右侧只要求在较小的 \(A\) 范围上接受类型实参，而左侧能接受所有
\(B\) 的子类型；在 \(X<:A\) 下又有 \(X<:B\)。核规则要求两边上界完全
相同，因而不能推出。把量词体的 \(X,B\) 换成其他满足协变关系的类型，还可
得到更多例子。
\taplsolutionbacklink{ex:26.2.2}{返回练习26.2.2}
\end{taplsolutionentrybox}

\begin{taplsolutionentrybox}
\phantomsection\label{sol:translator:26.3.1}
\textbf{26.3.1 解答}

展开序对：
\[
\forall X.(S_1\to S_2\to X)\to X
\quad\text{与}\quad
\forall X.(T_1\to T_2\to X)\to X.
\]
量词上界同为 \(\tapltype{Top}\)，可用核 S-All。量词体都是箭头；最外层
参数位置逆变，所以只需证明
\[
T_1\to T_2\to X<:S_1\to S_2\to X.
\]
再两次使用 S-Arrow，正好分别需要 \(S_1<:T_1\)、\(S_2<:T_2\)，结果位置
\(X<:X\) 由 S-Refl 得到。
\taplsolutionbacklink{ex:26.3.1}{返回练习26.3.1}
\end{taplsolutionentrybox}

\begin{taplsolutionentrybox}
\phantomsection\label{sol:translator:26.3.6}
\textbf{26.3.6 解答}

一般不能。子类型关系只能让一个已有类型忘掉信息，不能在项内部抽象出一个
可由调用者任意选择的类型参数。全称类型的关键能力是：同一值可在许多互不
相关的类型处实例化，并在结果类型中保持这些类型之间的联系。仅靠
\(\tapltype{Top}\)、箭头和 T-Sub 无法恢复这种依赖关系。
\taplsolutionbacklink{ex:26.3.6}{返回练习26.3.6}
\end{taplsolutionentrybox}

\begin{taplsolutionentrybox}
\phantomsection\label{sol:translator:26.4.16}
\textbf{26.4.16 解答}

完全系统的证明框架不变，但所有涉及 S-All 的辅助引理都必须同时处理变化的
上界。收窄和类型替换证明在量词情形分别使用完全 S-All 的第一个前提来调整
上界，再在较窄上下文中处理量词体。反演引理中，全称类型两边的上界不再相等，
而只满足完全规则给出的逆变关系。

把这些加强后的引理代入保持性证明，类型 beta 归约仍由类型替换引理处理；
进展性只关心闭值的外层形式，不依赖两个上界相等，因此典范形式和主证明保持
原有结构。由此完全 \(F_{<:}\) 同样满足保持性与进展性。
\taplsolutionbacklink{ex:26.4.16}{返回练习26.4.16}
\end{taplsolutionentrybox}
